\documentclass[12pt]{article}
\usepackage[margin=0.75in,top=0.8in,bottom=0.65in]{geometry}
\usepackage{lmodern}
\usepackage{microtype}
\usepackage{fancyhdr}
\usepackage{titlesec}
\usepackage{enumitem}
\usepackage{lastpage}
\usepackage[hidelinks]{hyperref}

\setlength{\parindent}{0pt}
\setlength{\parskip}{0.38em}
\setlength{\headheight}{26pt}
\titleformat{\section}{\large\bfseries\scshape}{}{0pt}{}
\titleformat{\subsection}{\normalsize\bfseries}{}{0pt}{}
\pagestyle{fancy}
\fancyhf{}
\fancyhead[C]{\footnotesize Lit Example Reader \quad Assignment ID: U1L4LE \quad Page \thepage\ of \pageref{LastPage}}
\renewcommand{\headrulewidth}{0.5pt}

\newenvironment{playpassage}{\begin{quote}\small\setlength{\parskip}{0.25em}}{\end{quote}}
\newcommand{\speaker}[1]{\textbf{#1.}}

\begin{document}

\begin{center}
{\LARGE\bfseries Week 4 Lit Example Reader}\par\vspace{0.05em}
{\large\scshape Julius Caesar: The Middle Term of Future Danger}\par\vspace{0.15em}
\rule{0.78\textwidth}{0.5pt}
\end{center}

\textbf{Purpose:} This reader gives Julius Caesar passages for applying Week 4's logic habit: map a practical syllogism, name the major term, minor term, and middle term, and test whether the middle term really connects the premises to the conclusion.

\textbf{What to watch for:} Brutus and Cassius do not merely express feelings about Caesar. They try to make political fear look like proof. The logical question is: \textit{what term is supposed to bridge the evidence to the conclusion?}

\section*{Before the Passages: The Week 4 Logic Tool}

A syllogism uses two premises to support a conclusion. In a practical literary argument, the exact wording may be messy, but the map still matters:
\begin{itemize}[leftmargin=1.5em,itemsep=0.2em]
\item \textbf{Major term:} the predicate of the conclusion.
\item \textbf{Minor term:} the subject of the conclusion.
\item \textbf{Middle term:} the term that appears in the premises and is supposed to connect the minor term to the major term.
\end{itemize}

The middle term is the pressure point. If it does not truly apply, the conclusion may only look proved.

\section*{Passage 1: Act 2, Scene 1 --- Brutus and the Serpent's Egg}

\textbf{Context:} Brutus is alone, thinking about whether Caesar should be killed before he becomes king. Caesar has not yet become a tyrant. Brutus has to reason from what Caesar may become.

\begin{playpassage}
\speaker{Brutus} It must be by his death: and for my part,\par
I know no personal cause to spurn at him,\par
But for the general. He would be crown'd:\par
How that might change his nature, there's the question.\par
It is the bright day that brings forth the adder;\par
And that craves wary walking. Crown him?---that;---\par
And then, I grant, we put a sting in him,\par
That at his will he may do danger with.\par
The abuse of greatness is, when it disjoins\par
Remorse from power: and, to speak truth of Caesar,\par
I have not known when his affections sway'd\par
More than his reason. But 'tis a common proof,\par
That lowliness is young ambition's ladder,\par
Whereto the climber-upward turns his face;\par
But when he once attains the upmost round,\par
He then unto the ladder turns his back,\par
Looks in the clouds, scorning the base degrees\par
By which he did ascend. So Caesar may.\par
Then, lest he may, prevent. And, since the quarrel\par
Will bear no colour for the thing he is,\par
Fashion it thus; that what he is, augmented,\par
Would run to these and these extremities:\par
And therefore think him as a serpent's egg\par
Which, hatch'd, would, as his kind, grow mischievous,\par
And kill him in the shell.
\end{playpassage}

\subsection*{Reader Notes}

Here is your most assisted example. Brutus's likely conclusion is that Caesar should be stopped before he is crowned. A possible middle term is something like \textit{one who will probably become a dangerous tyrant if given supreme power}. Notice how much weight that term carries. Brutus admits he has no personal cause against Caesar and has not known Caesar's passions to rule his reason; then he shifts to what ambition commonly does.

\textbf{Reading task:} Underline Brutus's conclusion. Box the lines where he admits what he has not proved about Caesar. Star the image that makes future danger feel vivid.

\newpage

\section*{Passage 2: Act 1, Scene 2 --- Cassius and the Colossus}

\textbf{Context:} Cassius wants Brutus to resist Caesar. His argument is less orderly than Brutus's orchard speech. Watch how he moves from Caesar's greatness, Roman honor, and Brutus's family name toward a political conclusion.

\begin{playpassage}
\speaker{Cassius} Why, man, he doth bestride the narrow world\par
Like a Colossus, and we petty men\par
Walk under his huge legs and peep about\par
To find ourselves dishonourable graves.\par
Men at some time are masters of their fates:\par
The fault, dear Brutus, is not in our stars,\par
But in ourselves, that we are underlings.\par
Brutus and Caesar: what should be in that ``Caesar''?\par
Why should that name be sounded more than yours?\par
Write them together, yours is as fair a name;\par
Sound them, it doth become the mouth as well;\par
Weigh them, it is as heavy; conjure with 'em,\par
Brutus will start a spirit as soon as Caesar.\par
Now, in the names of all the gods at once,\par
Upon what meat doth this our Caesar feed,\par
That he is grown so great? Age, thou art shamed!\par
Rome, thou hast lost the breed of noble bloods!\par
When could they say till now, that talk'd of Rome,\par
That her wide walls encompass'd but one man?
\end{playpassage}

\subsection*{Reader Notes}

Cassius gives you evidence to sort, not a finished diagram. He points to Caesar's public greatness, Brutus's equal name, Roman shame, and the danger of living under one man. Do not let the insult or imagery do the mapping for you. Decide which term is supposed to connect Caesar's position to the conclusion Cassius wants.

\textbf{Reading task:} Mark three phrases Cassius uses as evidence. In the margin, write one possible middle term, but do not decide yet whether it is proved.

\section*{How to Use These Passages on the Worksheet}

The worksheet asks you to map Shakespearean reasoning into premises, conclusion, and terms. Quote briefly if useful, but the answer should be a logic map, not a plot summary.

\textbf{Strong college-level answers} should name the conclusion, identify the middle term, and explain whether the middle term is proved strongly enough to carry the argument.

\end{document}
